\QHeading{第13题　在$\bigcirc$里填上“$>$”、“$<$”或“$=$”。}
\par\noindent


\PlainItem{2千克{\color{answerred}$>$}200克}
\par

\PlainItem{0.8吨{\color{answerred}$>$}80千克}
\PlainItem{3元5角{\color{answerred}$>$}3.05元}
\PlainItem{5平方分米{\color{answerred}$<$}0.5平方米}
\PlainItem{5.5分米{\color{answerred}$=$}55厘米}
\PlainItem{1.5吨{\color{answerred}$=$}1500千克}
\PlainItem{6.6平方分米{\color{answerred}$>$}66平方厘米}
\PlainItem{0.08米{\color{answerred}$=$}8厘米}
\par\noindent

\QHeading{第14题　解决问题。}

\Problem{① 声音在空气中的传播速度大约为每秒340米，合每分钟多少千米?}
\par


\MethodTitle{方法一}
\Answer{\Formula{\begin{aligned}[t]&\phantom{={}}340\times60\div1000\\&=20400\div1000\\&=20.4\text{（千米）}\end{aligned}}}
\Answer{答：合每分钟20.4千米。}

\MethodTitle{方法二：先将每秒米数换成千米}
\Answer{340米等于0.34千米，1分钟有60秒。}
\Answer{\Formula{\begin{aligned}[t]&\phantom{={}}0.34\times 60\\
&=20.4\text{（千米）}\end{aligned}}}
\Answer{答：合每分钟20.4千米。}
\Problem{② 100吨煤可炼焦炭95吨。照这样计算，10吨、1000吨煤分别可炼焦炭多少吨?}

\MethodTitle{方法一}
\Answer{10吨煤可炼焦炭：}
\Answer{\Formula{\begin{aligned}[t]&\phantom{={}}95\div(100\div10)\\&=95\div10\\&=9.5\text{（吨）}\end{aligned}}}
\Answer{1000吨煤可炼焦炭：}
\Answer{\Formula{\begin{aligned}[t]&\phantom{={}}95\times(1000\div100)\\&=95\times10\\&=950\text{（吨）}\end{aligned}}}
\Answer{答：10吨煤可炼焦炭9.5吨。}
\Answer{答：1000吨煤可炼焦炭950吨。}

\MethodTitle{方法二：先求1吨煤的产量，再放大}
\Answer{先用95除以100求1吨煤可炼焦炭多少吨。}
\Answer{\textbf{每吨煤可炼的焦炭}\par \Formula{\begin{aligned}[t]&\phantom{={}}95\div 100\\
&=0.95\text{（吨）}\end{aligned}}}
\Answer{\textbf{10吨煤可炼的焦炭}\par \Formula{\begin{aligned}[t]&\phantom{={}}0.95\times 10\\
&=9.5\text{（吨）}\end{aligned}}}
\Answer{\textbf{1000吨煤可炼的焦炭}\par \Formula{\begin{aligned}[t]&\phantom{={}}0.95\times 1000\\
&=950\text{（吨）}\end{aligned}}}
\Answer{答：10吨煤可炼焦炭9.5吨。\par 答：1000吨煤可炼焦炭950吨。}
\Problem{③ 100千克某种花生仁能榨出41千克花生油，1吨这样的花生仁能榨出多少千克花生油? 1千克这样的花生仁能榨出多少千克花生油?}

\MethodTitle{方法一}
\Answer{1吨花生仁的出油量：}
\Answer{\Formula{\begin{aligned}[t]&\phantom{={}}41\times(1000\div100)\\&=41\times10\\&=410\text{（千克）}\end{aligned}}}
\Answer{1千克花生仁的出油量：}
\Answer{\Formula{\begin{aligned}[t]&\phantom{={}}41\div100\\&=0.41\text{（千克）}\end{aligned}}}
\Answer{答：1千克花生仁能榨油0.41千克。}
\Answer{答：1吨花生仁能榨油410千克。}

\MethodTitle{方法二：先求1千克花生仁的出油量}
\Answer{1吨等于1000千克，先求1千克，再乘1000。}
\Answer{\textbf{1千克花生仁出油量}\par \Formula{\begin{aligned}[t]&\phantom{={}}41\div 100\\
&=0.41\text{（千克）}\end{aligned}}}
\Answer{\textbf{1吨花生仁出油量}\par \Formula{\begin{aligned}[t]&\phantom{={}}0.41\times 1000\\
&=410\text{（千克）}\end{aligned}}}
\Answer{答：1千克能榨出0.41千克花生油。\par 答：1吨能榨出410千克花生油。}
\Problem{④ 学校准备给一间会议室铺正方形地砖，每块正方形地砖的边长是60厘米，刚好每行铺20块，每列铺24块。这间会议室的面积是多少平方米?}

\MethodTitle{方法一}
\Answer{\Formula{\begin{aligned}[t]&\phantom{={}}(6\times6)\times(20\times24)\div100\\&=36\times480\div100\\&=17280\div100\\&=172.8\text{（平方米）}\end{aligned}}}
\Answer{$60$厘米$=6$分米。\par 答：会议室面积\par 是172.8平方米。}

\MethodTitle{方法二：先求会议室的长、宽}
\Answer{一行20块、一列24块，先按每块边长60厘米求总长、总宽。为了只用已学的整数乘法，先用分米计算。}
\Answer{\textbf{会议室的一边}\par \Formula{\begin{aligned}[t]&\phantom{={}}6\times 20\\
&=120\text{（分米）}\end{aligned}}}
\Answer{\textbf{会议室的另一边}\par \Formula{\begin{aligned}[t]&\phantom{={}}6\times 24\\
&=144\text{（分米）}\end{aligned}}}
\Answer{\textbf{面积并换成平方米}\par \Formula{\begin{aligned}[t]&\phantom{={}}120\times 144\div 100\\
&=17280\div 100\\
&=172.8\text{（平方米）}\end{aligned}}}
\Answer{答：会议室的面积是172.8平方米。}
